\section{Executable Manipulation Checks}
\label{sec:method}

\subsection{Giving \texorpdfstring{$f$}{f} empirical content}
\label{sec:empirical-content}

The definition of an information processing agent is that it executes a manipulation
$f(C_i, y, x)$ on $y$. There are no restrictions on $f$, meaning any post-hoc observation about a
model's behavior given a certain prompt can be claimed to be implementing an IPA, which makes the
definition unfalsifiable and purely descriptive. To address this, we will demand a mechanism to
check that the claimed manipulation was indeed performed.

\begin{definition}[executable manipulation check]
\label{def:check}
An executable manipulation check for a role $i$ is a scalar-valued function $m_i$, taking as input
structured output records, specified prior to running the experiment, where a higher value of $m_i$
indicates that the structured output records are consistent with the claimed manipulation performed
by $C_i$. For a role's performance to be ascribed to its operation, its manipulation check $m_i$
must discriminate between the structured output records of the role, and a length-matched control
condition.
\end{definition}

Throughout, \emph{implements} is operational: the role must produce the predeclared signature in its
structured records, and no claim is made about latent computation that leaves no such trace.

In experimental psychology, a manipulation check can confirm that the experimental manipulation was
conducted as planned, before making any conclusions about the experimental condition's effects
\citep{ejelov2020manipulation}. They're common in published studies: \citet{ejelov2020manipulation}
report their usage across 207 published experiments. In multi-agent LLM literature, they are rare;
\S\ref{sec:related} names the closest cases.

While we use manipulation checks, we also acknowledge that their use in psychology is itself
controversial, both because they can trigger self-reinforcing or self-defeating dynamic processes
\citep{hauser2018necessary} and because the majority of checks reported do not account for these
dynamics \citep{ejelov2020manipulation}. These concerns do not apply in this context, as our
manipulation checks are ex-post functions over already generated structured outputs, that only use
data produced by the benchmark. They are therefore non-interventionist, and do not `ripple' into the
experiment: the key concern with manipulation checks in human studies stems from a necessary human
intervention, which is not required when the check is executable code.

\begin{figure}[t]
\centering
\begin{tikzpicture}[
  font=\small,
  box/.style={draw, rounded corners=1pt, minimum height=6.5mm, inner xsep=4pt, align=center},
  gate/.style={draw, rounded corners=1pt, fill=blue!6, minimum height=6.5mm, inner xsep=4pt,
               align=center},
  arr/.style={-{Latex[length=1.6mm]}, gray!70},
]
  \node[box] (op)      {operation};
  \node[box, right=3mm of op]  (pe) {persona};
  \node[box, right=3mm of pe]  (ne) {neutral};

  \node[gray!70, font=\scriptsize\itshape, above=1mm of pe.north]
        {length-matched; identical output contract};
  \draw[gray!55] ([xshift=-1mm,yshift=1.5mm]op.north west) -- ([xshift=1mm,yshift=1.5mm]ne.north east);

  \node[box, below=8mm of pe] (model) {same model,\\fixed decoding};
  \node[box, right=5mm of model] (rec) {structured\\records};
  \node[box, right=5mm of rec]   (chk) {check $m_i$ vs.\\neutral: $d$, CI};
  \node[gate, right=5mm of chk]  (gate) {gate: $q \le 0.05$ \textbf{and} $d \ge 0.40$};

  \node[gray!70, font=\scriptsize\itshape, align=center, above=1.5mm of gate.north]
        {attribution\\licensed only if};

  \draw[arr] (op.south) |- ([yshift=3mm]model.north west) -- ([xshift=-3mm]model.north);
  \draw[arr] (pe.south) -- (model.north);
  \draw[arr] (ne.south) |- ([yshift=3mm]model.north east) -- ([xshift=3mm]model.north);
  \draw[arr] (model) -- (rec);
  \draw[arr] (rec)   -- (chk);
  \draw[arr] (chk)   -- (gate);
\end{tikzpicture}
\caption{The check pipeline instantiated each role via three system prompts sharing one output
contract (the predeclared processing instruction, the role label alone, and no instruction at all).
All three traverse an identical path, so the contrast at the end isolates the prompt, and the gate
decides whether a benchmark delta may be read as an effect of the operation at all.}
\label{fig:pipeline}
\end{figure}

\subsection{Three prompt variants}
\label{sec:variants}

We present three conditions: operation vs neutral, which is the key contrast, and a persona
condition that disentangles the role condition from the operation it represents. We programmatically
test that all control conditions are padded so their length is within 10\% of the operation prompt
(to ensure that only the operation is helpful, not longer token strings), that our persona
conditions provably do not contain any words from the operation, and assert that none of the prompts
(when rendered) contain any medical terms. In all cases, we use a standardized structured output
format (a single JSON block with answer, abstain, confidence, list of atomic statements with whether
they have evidence, counterevidence and failures), which is necessary to enable all conditions to be
evaluated by the same set of metrics (the costs of this approach are discussed in
\S\ref{app:concedes}).

\subsection{The six checks}
\label{sec:six-checks}

For each condition, we specify a manipulation check. These manipulation checks (rate of
counterclaims and flags, number of contradictions found (corrected for abstention), number of atomic
statements made, diversity $\times$ concentration of hypotheses, number of sourced and unsourced
alternatives proposed and negative expected calibration error (NECE)) were all explicitly specified,
prior to data collection, by condition, and are only properties of the transformation represented by
the condition rather than the ``correctness'' of a solution (see Appendix~\ref{app:checks} for
details on how these checks are calculated).

\subsection{Decision rule and gate}
\label{sec:gate}

The checks are paired at the item level (the same item under all three variants) and per-item
differences are bootstrapped over 2,000 resamples. Both conditions in Figure~\ref{fig:pipeline} gate
on the number of observations, $n = 200$ --- a tight interval around a negligible difference will
pass a significance-only rule, and a large standardized effect with an interval spanning zero will
not offer evidence for a difference. The threshold is on a standardized but not absolute scale
because the six checks have no common unit (two return unbounded counts, one a rate difference in
$[-1, 1]$, one a distinct-item count). A single absolute threshold would be uninterpretable. Bounded
rates are sensitive to their inputs; rounding errors can seem like large moves.

The check runner exits non-zero if any of the roles fail to pass their checks, making them a gate on
a collection run instead of a report that will be produced and not read. If a role fails, it means
that this role hasn't validated the predeclared behavioural signature of its operation ---
equivalently, no matter how good it scored on the benchmark, this delta isn't licensed as evidence
for an effect of this operation.
